% SPDX-License-Identifier: CC-BY-NC-ND-4.0
% !TEX root = main.tex

\section{アーベル圏}

\begin{definition} \label{def: IsPreadditive}
局所小圏$\mathscr{C}$は次の条件を満たすとき\term{ぜんかほうけん}{前加法圏}{preadditive category}であるという.
\begin{itemize}
    \item 任意の$A, B \in \Ob(\mathscr{C})$に対して写像$\sigma_{(A,B)} \colon \Mor_{\mathscr{C}}(A, B) \times \Mor_{\mathscr{C}}(A, B) \to \Mor_{\mathscr{C}}(A, B)$が定まり次の条件を満たす.\
    ただし任意の$f, g \in \Mor_{\mathscr{C}}(A, B)$に対して$f + g \coloneq \sigma_{(A,B)}(f,g)$とする.
    \begin{itemize}
    \item ある$0_{A,B}^{+} \in \Mor_{\mathscr{C}}(A,B)$が存在して$(\Mor_{\mathscr{C}}(A,B), \sigma_{(A,B)}, 0_{A,B} )$はアーベル群.
    \item 任意の$A, B, C \in \Ob(\mathscr{C})$と$f,g \in \Mor_{\mathscr{C}}(A, B),\ h \in \Mor_{\mathscr{C}}(B, C)$に対して$h \circ (f + g) = h \circ f + f \circ g$成り立つ.
    \item 任意の$A, B, C \in \Ob(\mathscr{C})$と$f,g \in \Mor_{\mathscr{C}}(B, C),\ h \in \Mor_{\mathscr{C}}(A, B)$に対して$(f + g) \circ h = f \circ h + g \circ h$成り立つ.
    \end{itemize}
\end{itemize}
\end{definition}

\begin{definition} \label{def: zeroMorphism}
$\mathscr{C}$を零対象$Z$を持つ局所小圏とし$A, B \in \Ob({\mathscr{C}})$とする.\
$\iexcl_B \in \Mor_{\mathscr{C}}(Z, B)$と$!_A \in \Mor_{\mathscr{C}}(A, Z)$をそれぞれただ1つ存在する射とする.\
このとき$0_{A,B} \coloneq !_A \circ \iexcl_B$を\term{れいしゃ}{零射}{zero morphism}という.
\end{definition}

\begin{lemma} \label{lem: additive_zero_comp}
$\mathscr{C}$を零対象$Z$を持つ前加法圏とする.\
任意の$A, B, C \in \Ob(\mathscr{C})$と$g \in \Mor_{\mathscr{C}}(B, C),\ h \in \Mor_{\mathscr{C}}(C, A)$に対して次の図式が可換になる.
% https://q.uiver.app/#q=WzAsNixbMCwwLCJCIl0sWzEsMCwiQSJdLFsxLDEsIkMiXSxbMywwLCJCIl0sWzQsMCwiQSJdLFs0LDEsIkMiXSxbMSwwLCIwX3tBLEJ9XnsrfSIsMl0sWzAsMiwiZyIsMl0sWzEsMiwiMF97QSxDfV57K30iXSxbNSwzLCIwX3tDLEJ9XnsrfSJdLFs0LDMsIjBfe0EsIEJ9XnsrfSIsMl0sWzUsNCwiaCIsMl1d
\[\begin{tikzcd}
	B & A && B & A \\
	& C &&& C
	\arrow["g"', from=1-1, to=2-2]
	\arrow["{0_{A,B}^{+}}"', from=1-2, to=1-1]
	\arrow["{0_{A,C}^{+}}", from=1-2, to=2-2]
	\arrow["{0_{A, B}^{+}}"', from=1-5, to=1-4]
	\arrow["{0_{C,B}^{+}}", from=2-5, to=1-4]
	\arrow["h"', from=2-5, to=1-5]
\end{tikzcd}\]
\end{lemma}

\begin{proof}
$g \circ 0_{A,B}^{+} = g \circ (0_{A,B}^{+} + 0_{A,B}^{+}) = g \circ 0_{A,B}^{+} + g \circ 0_{A,B}^{+}$であり両辺$g \circ 0_{A,B}^{+}$を引けば$0_{A,C}^{+} = g \circ 0_{A,B}^{+}$となる.\
同様に$0_{A,B}^{+} \circ h = 0_{C,B}^{+}$となる.
\end{proof}

\begin{proposition} \label{prop: additive_zero_eq_zero_morphism}
$\mathscr{C}$を零対象$Z$を持つ前加法圏とする.\
任意の$A,B \in \Ob(\mathscr{C})$に対して次が成り立つ.
\begin{center}
    $0_{A,B}^{+} = 0_{A,B}$
\end{center}
\end{proposition}

\begin{proof}
$\Mor_{\mathscr{C}}(Z, B)$はシングルトンなので$\iexcl_B = 0_{Z,B}^{+}$である.\
\cref{lem: additive_zero_comp}により$0_{A,B} = !_A \circ 0_{Z,B}^{+} = 0_{A,B}^{+}$となる.
\end{proof}

\begin{definition} \label{def: WalkingPairCategory}
圏$\mathbf{2}$を$\Ob(\mathbf{2}) \coloneq 2,\ \Mor(\mathbf{2}) \coloneq \{\id_0, \id_1\}$として定める.
\end{definition}

\begin{definition} \label{def: pairDiagramFunctor}
$\mathscr{C}$を局所小圏とし$A,B \in \Ob(\mathscr{C})$とする.\
共変関手$D_{A,B} \colon \mathbf{2} \to \mathscr{C}$を次のようにして定める.
\begin{itemize}
    \item $D_{A,B}(0) \coloneq A,\ D_{A,B}(1) \coloneq B$とする.
    \item 任意の$j \in 2$に対して$D_{A,B}(\id_j) \coloneq \id_{D_{A,B}(j)}$とする.
\end{itemize}
\end{definition}

\begin{definition} \label{def: binaryProduct}
$\mathscr{C}$を局所小圏とし$A,B \in \Ob(\mathscr{C})$とする.\
$(L,(L, \pi))$を$D_{A,B}$の極限とする.\
このとき$A \times B \coloneq L$を\term{にこうせき}{2項積}{binary product}といい$\pi_0, \pi_1$を\term{しゃえい}{射影}{projection}という.
\end{definition}

\begin{definition} \label{def: IsAdditive}
前加法圏$\mathscr{C}$は次の条件を満たすとき\term{かほうけん}{加法圏}{additive category}であるという.
\begin{itemize}
    \item 零対象$Z \in \Ob(\mathscr{C})$が存在.
    \item 任意の$A, B \in \Ob(\mathscr{C})$に対して$D_{A,B} \colon \mathbf{2} \to \mathscr{C}$の極限$(A \times B, (A \times B, \pi))$が存在.
\end{itemize}
\end{definition}

\begin{definition} \label{def: WalkingParallelPairCategory}
圏$\mathbf{P}$を$\Ob(\mathbf{P}) \coloneq 2$とし$\Mor{\mathscr{C}} \coloneq \{\id_0, \id_1, u, v\}$として定める.\
ただし$|\Mor_{\mathbf{P}}(0, 1)| \coloneq 2$である.\
これを\term{へいこうついのけん}{平行対の圏}{walking parallel pair category}という.\
平行対の圏を$\bullet \rightrightarrows \bullet$と書くこともある.
% https://q.uiver.app/#q=WzAsMixbMCwwLCIwIl0sWzEsMCwiMSJdLFswLDEsInUiLDAseyJjdXJ2ZSI6LTF9XSxbMCwwLCJcXG1hdGhybXtpZH1fMCIsMCx7ImFuZ2xlIjotOTB9XSxbMSwxLCJcXG1hdGhybXtpZH1fMSIsMCx7ImFuZ2xlIjo5MH1dLFswLDEsInYiLDIseyJjdXJ2ZSI6MX1dXQ==
\[\begin{tikzcd}
	0 & 1
	\arrow["{\mathrm{id}_0}", from=1-1, to=1-1, loop, in=145, out=215, distance=10mm]
	\arrow["u", curve={height=-6pt}, from=1-1, to=1-2]
	\arrow["v"', curve={height=6pt}, from=1-1, to=1-2]
	\arrow["{\mathrm{id}_1}", from=1-2, to=1-2, loop, in=325, out=35, distance=10mm]
\end{tikzcd}\]
\end{definition}

\begin{definition} \label{def: kernel}
$\mathscr{C}$を加法圏とし$f \in \Mor({\mathscr{C}})$とする.\
$\mathbf{P}$を平行対の圏とし共変関手$K_f \colon \mathbf{P} \to \mathscr{C}$を次のようにして定める.
\begin{itemize}
    \item $K_f(0) \coloneq \dom(f),\ K_f(1) \coloneq \cod(f)$とする.
    \item $K_f(u) \coloneq f,\ K_f(v) \coloneq 0_{\dom(f), \cod(f)}$とする.
\end{itemize}
$K_f$の極限$(L_{K_f}, (L_{K_f}, \pi_{K_f}))$を$f$の\term{かく}{核}{kernel}といい$\ker(f) \coloneq L_{K_f},\ k_f \coloneq (\pi_{K_f})_0 \colon \ker(f) \to \dom(f)$とする.\
つまり任意の$X \in \Ob(\mathscr{C})$と$(X, \eta) \in \Cone(X,K_f)$に対して次の図式を可換にする$g \colon X \to \ker(f)$がただ1つ存在する.
% https://q.uiver.app/#q=WzAsNCxbMSwxLCJcXG1hdGhybXtLZXJ9KGYpIl0sWzAsMiwiXFxtYXRocm17ZG9tfShmKSJdLFsyLDIsIlxcbWF0aHJte2NvZH0oZikiXSxbMSwwLCJYIl0sWzAsMSwia19mIiwyLHsibGFiZWxfcG9zaXRpb24iOjQwfV0sWzEsMiwiZiIsMCx7Im9mZnNldCI6LTF9XSxbMCwyLCIoXFxwaV97S19mfSlfMSIsMCx7ImxhYmVsX3Bvc2l0aW9uIjozMH1dLFsxLDIsIjBfe1xcbWF0aHJte2RvbX0oZiksXFxtYXRocm17Y29kfShmKX0iLDIseyJvZmZzZXQiOjF9XSxbMywxLCJcXGV0YV8wIiwyLHsiY3VydmUiOjN9XSxbMywyLCJcXGV0YV8xIiwwLHsiY3VydmUiOi0zfV0sWzMsMCwiZyIsMl1d
\[\begin{tikzcd}
	& X & \\
	& {\mathrm{Ker}(f)} \\
	{\mathrm{dom}(f)} && {\mathrm{cod}(f)}
	\arrow["g"', from=1-2, to=2-2]
	\arrow["{\eta_0}"', curve={height=18pt}, from=1-2, to=3-1]
	\arrow["{\eta_1}", curve={height=-18pt}, from=1-2, to=3-3]
	\arrow["{k_f}"'{pos=0.4}, from=2-2, to=3-1]
	\arrow["{(\pi_{K_f})_1}"{pos=0.3}, from=2-2, to=3-3]
	\arrow["f", shift left, from=3-1, to=3-3]
	\arrow["{0_{\mathrm{dom}(f),\mathrm{cod}(f)}}"', shift right, from=3-1, to=3-3]
\end{tikzcd}\]
\end{definition}

\begin{remark} \label{rem: kernel_cone_zero_leg}
\cref{def: kernel}の仮定のもと任意の$X \in \Ob(\mathscr{C})$と$(X, \eta) \in \Cone(X,K_f)$に対して$\eta_1 = 0_{\dom(f), \cod(f)} \circ f = 0_{X, \cod(f)}$である.
\end{remark}

\begin{definition} \label{def: cokernel}
$K_f$の余極限$(Q_{K_f}, (Q_{K_f}, \iota_{K_f}))$を$f$の\term{よかく}{余核}{cokernel}といい$\coker(f) \coloneq Q_{K_f},\ c_f \coloneq (\iota_{K_f})_0 \colon \cod(f) \to \coker(f)$とする.
\end{definition}

\begin{definition} \label{def: IsMonomorphism}
$\mathscr{C}$を圏とし$f \in \Mor(\mathscr{C})$とする.\
任意の$X \in \Ob(\mathscr{C})$と$g, h \in \Mor_{\mathscr{C}}(X, \dom(f))$に対して次の条件が成り立つとき$f$は\term{ものしゃ}{モノ射}{monomorphism}であるといい$f \colon \dom(f) \hookrightarrow \cod(f)$と書く.
\begin{center}
    $f \circ g = f \circ h \implies g = h$
\end{center}
\end{definition}

\begin{definition} \label{def: IsEpimorphism}
$\mathscr{C}$を圏とし$f \in \Mor(\mathscr{C})$とする.\
任意の$Y \in \Ob(\mathscr{C})$と$g, h \in \Mor_{\mathscr{C}}(\cod(f), Y)$に対して次の条件が成り立つとき$f$は\term{えぴしゃ}{エピ射}{epimorphism}であるといい$f \colon \dom(f) \twoheadrightarrow \cod(f)$と書く.
\begin{center}
    $g \circ f = h \circ f \implies g = h$
\end{center}
\end{definition}

\begin{definition} \label{def: IsAbelian}
加法圏$\mathcal{A}$は次の条件を満たすとき\term{あーべるけん}{アーベル圏}{Abelian category}であるという.
\begin{itemize}
    \item 任意の$f \in \Mor(\mathrm{C})$に対して$f$の核と余核が存在する.
    \item 任意のモノ射$f \in \Mor({\mathscr{C}})$に対してある$g \in \Mor_{\mathscr{C}}$が存在して$f = k_g$となる.
    \item 任意のエピ射$f \in \Mor({\mathscr{C}})$に対してある$g \in \Mor_{\mathscr{C}}$が存在して$f = c_g$となる.
\end{itemize}
\end{definition}

\begin{definition} \label{def: image}
$\mathscr{C}$を核と余核をもつ加法圏とし$f \in \Mor({\mathscr{C}})$とする.\
$\im(f) \coloneq \ker(c_f)$とし$i_f \coloneq k_{c_f}$とする.\
組$(\im(f), i_f)$を$f$の\term{ぞう}{像}{image}と呼ぶ.
\end{definition}